\begin{abstract}
  This paper presents a work-in-progress mechanization of the proof theory of
  the semi-axiomatic sequent calculus (SAX) in the proof assistant Beluga. SAX
  replaces all noninvertible inference rules with axioms, causing standard cut
  elimination to fail. Instead, SAX satisfies a modified cut elimination theorem
  permitting certain analytic cuts called snips to remain, where a snip requires
  the cut formula to be 'eligible', i.e. a subformula of the principal formula
  of an axiom. A central challenge in mechanizing SAX is therefore tracking
  eligibility throughout the proof theory. We address this by representing
  eligibility as a well-formedness predicate embedded directly within proof
  judgements rather than imposed through a separate context discipline, adapting
  a technique due to Crary for encoding substructural constraints as judgements
  over proof terms. This is analogous to how linearity conditions have been
  previously handled in mechanizations of session-typed systems. Using this
  technique, we present mechanized proofs of cut admissibility and the
  subformula property for a representative propositional fragment, and outline a
  roadmap for extending the mechanization to the full calculus. To our
  knowledge, this is the first mechanization of SAX proof theory.
\end{abstract}


% \begin{abstract}
%   % FIXME mention original authors in the abstract?
%   The semi-axiomatic sequent calculus (SAX), introduced by DeYoung, Pfenning,
%   and Pruiksma, arises from the standard sequent calculus $G_3$ by replacing all
%   noninvertible rules with axioms while leaving invertible rules unchanged. The
%   resulting proof system is equivalent to $G_3$ except that standard
%   cut-elimination fails, and is replaced by a modified form in which certain
%   analytic cuts, called snips, are permitted to remain in normal-form proofs
%   while preserving the subformula property. SAX admits a natural computational
%   interpretation as a form of shared memory concurrency.

%   % FIXME not sure if we should describe eligibility in the abstract
%   This paper presents a mechanization of the proof theory of SAX in the proof
%   assistant Beluga. A central challenge is the need to track eligibility, a
%   property of formula occurrences that identifies subformulas of principal
%   formulas introduced in axioms that are eligible for a snip. We describe how
%   eligibility is represented in Beluga as a predicate over derivations, adapting
%   a technique due to Crary for representing substructural constraints as
%   judgements over proof terms. We present mechanized proofs of the admissibility
%   of cut and the subformula property for a representative subset of SAX. To our
%   knowledge, this is the first mechanization of SAX proof theory.
% \end{abstract}
